\section{Games and logical interpretations}
\label{sec:interpretations}

\subsection{Full policies and endpoint plans}
\label{sec:policies}

For a public script, a full pure policy for role $i$ supplies an action at
every choice history owned by $i$, and a return value at each terminal
history. Each history includes the public parameter values and all earlier
choices, with their occurrence identities. The domains are well-typed
histories, including histories made unreachable by that role's own earlier
actions. A profile of such policies determines a history and a result record
by ordinary left-to-right evaluation.

This interpretation is set-theoretic. A general policy need not be definable
by a term of $\D$. Denote by $\llbracket\Local{i}{G}\rrbracket$ the analogous
set of dependent sums and functions in a set interpretation of the public
data. This notation does not assert that every set-theoretic endpoint has a
syntactic implementation.

Given a policy $\sigma_i$ and a history $h$, form its reduced endpoint plan
$\Plan_i(G,h,\sigma_i)$ recursively. At an owned choice, use the value
prescribed by $\sigma_i$ and retain its continuation for that value. At a
received choice, retain the function sending each possible value to its
corresponding continuation plan. At a terminal, return the policy's value.

\begin{theorem}[Execution adequacy for policy reduction]
\label{thm:adequacy}
For every full policy profile, set-theoretic execution of its reduced
endpoint plans has exactly the trace and return record of ordinary policy
evaluation.
\end{theorem}
\begin{proof}
At a terminal both evaluations select the same returns. At an owned choice,
the owner's reduced plan supplies the value prescribed by its policy at
the current history. Every other reduced plan is a function of that same
value, so all continuations are the reductions of the original policies at
the extended history. Apply the induction hypothesis to the continuation.
\end{proof}

The theorem preserves complete execution. It neither requires a
winning-strategy interpretation nor identifies policy tables with reduced
plans. Extending the comparison to an equivalence of strategic models must
state which unreachable histories are identified and how deviations and
randomization are represented.

\subsection{Objectives observe executions}

An objective can be any specified predicate $W_i(e)$ on $e:\Exec(G)$; a utility
can be a function $u_i(e)$ into an ordered payoff domain. An outcome tuple
$\langle[W_i(e)]\rangle_i$ is therefore one readout of the complete execution.
No exclusivity or agreement condition is imposed on the $W_i$.

For example, after (1), Alice may want the selected entry to be true, Bob may
want it to be false, and Carol may want the first position. These predicates
can conflict while the interaction remains well formed and executable.
Taking $R_i=\Unit$ does not certify any of them. Replacing $R_i$ by a
particular evidence type makes that proof a requirement of endpoint $i$;
it changes the interface explicitly.

Nash equilibrium, joint realizability, and adversarial guarantees can be
defined on the resulting policy game. Their quantification is additional
analysis data. It does not replace the profile space or complete evaluation
with a distinguished vector of winners.

\subsection{Classical quantifiers as a two-role interpretation}

The conventional evaluation game of first-order logic is recovered by fixing
a viewpoint and complementary objectives. We state this standard
specialization over finite nonempty structures, so it involves no choice of
noncomputable strategies over infinite domains.

Fix distinct roles $i,j$ and a prenex first-order sentence
$Q_1x_1\cdots Q_nx_n.\theta$, with quantifier-free matrix $\theta$.
Translate each existential binder to $\Q_i$, each universal binder to
$\Q_j$, and let all terminal result types be $\Unit$. Give role $i$ objective
$\theta$ and role $j$ objective $\neg\theta$. The positive objective belongs
to the selected role $i$, whichever role name is chosen.

\begin{proposition}[Classical evaluation-game correspondence]
In a finite nonempty first-order structure, the sentence is true iff role
$i$ has a policy that achieves $\theta$ against every policy of $j$.
\end{proposition}
\begin{proof}
Induct on the quantifier prefix. At an existential binder, truth requires
some value with a true continuation, and role $i$ can select that value.
At a universal binder, every value must have a true continuation; role $i$
can use the corresponding contingent continuation after observing $j$'s
choice. Finite branching permits all these choices to be assembled into
one policy. Conversely, a guaranteeing policy gives the existential witness
at an owned binder and a guaranteeing continuation after every opposing
choice. At the terminal the statement is satisfaction of $\theta$.
\end{proof}

This interpretation explains the existential/universal pair without assigning
permanent logical polarities to the roles of the global language. In a
general $K$-role interaction there is no distinguished opposing role.

\begin{proposition}[Role permutation is not general Boolean complement]
For $K>2$, no permutation of individual roles maps every one-hot outcome to
its coordinatewise Boolean complement.
\end{proposition}
\begin{proof}
A permutation preserves the number of nonzero coordinates. A one-hot vector
has one, while its complement has $K-1$.
\end{proof}

At $K=2$ these counts coincide, and swapping the complementary evaluation-game
roles has the familiar negation reading. For larger role sets, permutation
remains a symmetry of interaction. Logical complement of an objective remains
an operation in the objective language. The two operations have different
types and different laws.
